\section{Introduction}\label{sec:introduction}

Weak reversibility requires every reaction in a chemical reaction network
to have a directed return path. It is natural to ask whether this graph
condition alone prevents concentrations from approaching zero or infinity.
Permanence asks for more: after a transient, all trajectories allowed by
the same conservation laws must satisfy the same positive lower and finite
upper bounds. We prove this conclusion for arbitrary finite weakly
reversible networks with fixed positive reaction rates.

\subsection{The problem and main result}

Let $d\geq1$. A finite reaction network consists of a set of
\emph{complexes} $C\subset\Z_{\geq0}^d$ and a set of directed
\emph{reactions} $R\subset C\times C$. We write $y\to y'$ for
$(y,y')\in R$. The network is \emph{weakly reversible} if for every
$y\to y'\in R$ there is a directed path from $y'$ to $y$. A
\emph{linkage class} is a connected component of the underlying undirected
graph. Assign a positive constant $k_{y\to y'}$ to each reaction. The
mass-action equation is
\begin{equation}\label{eq:mass-action}
  \dot x=f(x):=\sum_{y\to y'\in R} k_{y\to y'}x^y(y'-y),
  \qquad x^y:=\prod_{i=1}^d x_i^{y_i}.
\end{equation}
The \emph{stoichiometric subspace} and the \emph{positive stoichiometric
class} of $c\in\R_{>0}^d$ are
\[
  S:=\operatorname{span}_{\R}\{y'-y:y\to y'\in R\},
  \qquad P:=(c+S)\cap\R_{>0}^d.
\]
Since $f(x)\in S$, any positive solution starting in $P$ remains in $P$
for as long as it exists. The class $P$ need not be bounded.

A global positive trajectory is \emph{persistent} if
$\liminf_{t\to\infty}x_i(t)>0$ for every $i$. These lower bounds may
depend on the trajectory. We call the system \emph{permanent on $P$} if
every solution starting in $P$ is global and positive, and there is one
compact set $K\subset P$ that eventually contains each such solution.
The entry time may depend on the initial point, but $K$ may not. This is
the classwise uniformity in the permanence question
\cite[Definition~4.1]{BorosHofbauer2020}.

\begin{theorem}\label{thm:main}
Let $C\subset\Z_{\geq0}^d$ be finite, let $R\subset C\times C$ be
weakly reversible, and fix positive rate constants in
\eqref{eq:mass-action}. For every positive stoichiometric class $P$ there
is a nonempty compact convex set $K_P\subset P$ with the following
properties. Every initial point $x^0\in P$ gives a unique global positive
solution, and there is a finite $T(x^0)\geq0$ such that
\[
  x(t)\in K_P\qquad(t\geq T(x^0)).
\]
Moreover, $K_P$ is forward invariant. In particular, there is
$\eps_P\in(0,1)$, depending only on the network, the fixed rates, and
$P$, such that every such solution satisfies
\begin{equation}\label{eq:uniform-bounds}
  \eps_P\leq x_i(t)\leq\eps_P^{-1}
  \qquad(1\leq i\leq d,\ t\geq T(x^0)).
\end{equation}
\end{theorem}

Thus the fixed-positive-rate, weakly reversible formulation of the
permanence conjecture holds, including on unbounded classes. The
uniformity in \eqref{eq:uniform-bounds} is within each fixed class and
rate choice; the transient time is allowed to depend on $x^0$.

\subsection{Historical context and the contribution}

The structural theory of mass-action systems begins with the work of
Horn and Jackson and of Feinberg
\cite{HornJackson1972,Feinberg1972}. A system is \emph{complex-balanced}
if it has a positive state at which the total incoming and outgoing
reaction fluxes agree at each complex. For a complex-balanced system,
the classical theory gives a unique positive equilibrium in each
positive stoichiometric class and
an entropy-type Lyapunov function. Complex balance imposes conditions
on the rates as well as the graph; weak reversibility by itself does not
impose these balance identities. Feinberg's discussion of weakly
reversible systems \cite[Remark~6.1.E]{Feinberg1987} raised the question
of excluding limiting states with a zero concentration. Persistence
excludes approach to the boundary even along subsequences, and
permanence adds upper bounds and uniformity over initial points.

Important partial results developed both the dynamical picture and the
geometric methods. Anderson proved that every positive trajectory of a
complex-balanced system with one linkage class converges to its class's
positive equilibrium
\cite[Corollary~4.11]{Anderson2011}. Craciun, Nazarov, and Pantea proved
permanence for two-species weakly reversible systems, as part of their
more general theorem for endotactic networks
\cite[Theorem~6.4]{CraciunNazarovPantea2013}. Their endotactic condition
expresses an inward-pointing property of the reactions seen in every
direction. Pantea proved persistence of bounded
weakly reversible trajectories when $\dim S=2$, with no restriction to
two species \cite[Theorem~5.1]{Pantea2012}. Gopalkrishnan, Miller, and Shiu proved
permanence under the geometric condition of strong endotacticity,
which includes every weakly reversible network with one linkage class
\cite[Theorems~1.1 and~1.3]{GopalkrishnanMillerShiu2014}.
Boros and Hofbauer later gave another proof of the single-linkage
permanence theorem, producing a compact forward-invariant absorbing
set \cite[Theorem~4.2]{BorosHofbauer2020}. These permanence results
also permit rates varying in time within fixed positive bounds.

Craciun's toric-differential-inclusion approach constructs invariant
regions by comparing monomials in logarithmic coordinates. The revised
manuscript of September~23, 2026 proves persistence for bounded positive
trajectories of weakly reversible systems, including rates bounded above
and away from zero; its complex-balanced application establishes global
convergence \cite[Section~8, final remark]{Craciun2026}. The general
weakly reversible persistence statement there assumes boundedness of
the trajectory. Our theorem treats fixed rates and supplies one
eventual bound for every trajectory in a class.

The immediate input to the present argument is the affine approximation
lemma and trapping construction of \cite{OpenAI2026}. That paper proves
that every prescribed positive initial point belongs to a compact
convex forward-invariant set contained in the positive orthant. Such sets
give boundedness and persistence, but may depend on the initial point.
We reproduce the affine lemma and
its proof in Section~\ref{sec:affine}, and adapt the trapping construction
in Section~\ref{sec:trapping}. The additional arguments establish strict
increase away from a trapping set and connect trapping sets at different
scales. They turn bounds for individual trajectories into permanence
on a fixed class. Monomial comparisons and invariant regions have the
earlier antecedents just described; the affine approximation lemma is
credited to \cite[Lemma~2.1 and Corollary~2.2]{OpenAI2026}.

\subsection{Proof overview}

Fix one representative $c\in P$. For $0<h<1$, set
$B_h=[h,h^{-1}]^d$. Every $x\in B_h$ has a unique representation
$x=(h^{p_1},\ldots,h^{p_d})$ with $p\in[-1,1]^d$. The affine lemma
constructs a finite concave minimum
\[
  F_h(x)=\min_{\lambda\in\Lambda}
     \bigl(r_\lambda\cdot x+b_\lambda(h)\bigr),
\]
whose slopes $r_\lambda\in[-1,1]^d$ are independent of $h$. A label
$\lambda$ is \emph{active} at $x$ when it attains this minimum.
For any prescribed tolerance $E:[-1,1]^d\to(0,\infty)$, the construction
ensures
$\norm{r_\lambda-p}_\infty<E(r_\lambda)$ for every active label,
uniformly for small $h$. We choose the tolerance so that every nonzero
comparison $r_\lambda\cdot v$, where $v$ is a coordinate vector or a
difference of distinct complexes, has the same sign as $p\cdot v$.
To preserve these signs, our tolerance shrinks near a hyperplane
$r\cdot v=0$; on the hyperplane, the zero comparison imposes no
restriction. This is why the affine lemma allows a discontinuous
tolerance. When $r_\lambda\cdot v\ne0$, finiteness of the slope family
also gives a uniform positive lower bound for $\abs{p\cdot v}$.

At sufficiently small scales, the labels active at $c$ have zero slope,
so $F_h(c)$ is the global maximum of $F_h$. Its plateau
\[
  K_h:=\{x\in\R^d:F_h(x)=F_h(c)\}
\]
is compact and convex, contains $c$, and satisfies
$K_h\subset\operatorname{int}B_{\sqrt h}$
(Proposition~\ref{prop:plateau}). The restriction to $P$ now supplies
the decisive strictness condition. A member of the affine family whose slope lies
in $S^\perp$ is constant on $c+S$, so its value on $P$ is at least
$F_h(c)$. Such a label cannot be active at a point of $P\setminus K_h$.

For any active label $\lambda$, order the complexes by the levels
$r_\lambda\cdot y$, and separate consecutive levels by cuts. If a
downward reaction crosses a cut, its directed return path contains an
upward crossing. Monomials at lower source levels dominate those at
higher levels, so one upward flux exceeds the entire downward flux for
sufficiently small $h$. The net upward flux across each cut is therefore
nonnegative, and is positive whenever a reaction crosses it. Summing
these fluxes times the gaps between levels gives
$r_\lambda\cdot f(x)\geq0$, with strict inequality when
$r_\lambda\notin S^\perp$. Compactness at each fixed scale makes the
strict inequality uniform over all active labels to which it applies
(Proposition~\ref{prop:increase}). Thus $F_h$ is nondecreasing along
solution segments in $B_h$ and increases at a fixed positive rate on
$B_h\cap(P\setminus K_h)$. A trajectory in $P$ that stays in $B_h$
for all later times must reach $K_h$ in finite time and remain there.

Every initial point in $P$ belongs to some $K_h$, giving global
existence first. The inclusion $K_h\subset B_{\sqrt h}$ then permits
passage to a larger scale. Repeated square roots, capped at a scale
$h_P$ fixed by the class, reach $h_P$ after finitely many steps.
Section~\ref{sec:permanence} carries out this iteration and obtains
$K_P=K_{h_P}\cap(c+S)$, with $h_P$ chosen independently of the initial
point.
